\documentclass[10pt,a4paper]{journal}
\usepackage{adjustbox,amsmath,amsfonts,bbm,caption,float,geometry,graphicx,hyperref}
\usepackage{cleveref,plex-serif,plex-sans,plex-mono,pdfpages}
\usepackage{listings,mathrsfs,mdwlist,pgfplots,siunitx,tikzit,wrapfig}
\geometry{verbose,tmargin=2cm,bmargin=3cm,lmargin=3cm,rmargin=3cm}
%\usepackage[dvips,xetex]{graphicx}
%\usepackage{ifpdf,mla}
\DeclareMathOperator{\Tr}{Tr}
\DeclareMathOperator{\re}{Re}
\DeclareMathOperator{\im}{Im}
\DeclareMathOperator{\ID}{\mathbbm{1}}
\usepackage[normalem]{ulem}
\usepackage[greek,irish,british]{babel}

%Cleveref config
%\crefrangeformat{equation}{(#3#1#4–#5#2#6)}
\crefname{equation}{equation}{equations}
\crefname{figure}{figure}{figures}
\crefname{table}{table}{tables}
%Allow aligned equations to be split across pages
\allowdisplaybreaks

%\usepackage[irish]{isodate}
%\babelprovide[import=el, onchar=ids fonts]{greek}
%\pagestyle{plain}
%
%Glossary setup
%\usepackage[automake,acronym,symbols,nomain,nogroupskip,nonumberlist,sort=use,record]{glossaries-extra}
%\makeglossaries
%\loadglsentries{./Glossary.tex}
%Index setup
%\usepackage{showidx}

\usetikzlibrary{shapes, arrows, fit}
\input{Tikz/Arrows.tikzstyles}

%Removing all references to Bastardised English
\newcommand{\colour}{\color}
\newcommand{\textcolour}{\textcolor}
\newcommand{\centring}{\centering}
\newenvironment{itemise}{\begin{itemize}}{\end{itemize}}
\newenvironment{centre}{\begin{center}}{\end{center}}

%Source Code
\lstset{language=C,basicstyle=\ttfamily\footnotesize,breaklines=true,numbers=left,tabsize=3,extendedchars=true,
		inputencoding=utf-8,commentstyle=\ttfamily\colour{orange},keywordstyle=\colour{blue}}
\title{The Clover Term in QC$_2$D}
\shorttitle{The Clover Term in QC$_2$D}
\subtitle{Or potentially a novel rapid method of developing a permenant brain injury.}
\author{(Still not a) Dr. Dale Lawlor}
\shortauthor{Lawlor}
\begin{document}
	\maketitle
	This is a brief derivation of the clover action for an $SU(2)$ HMC code at non-zero chemical potential (QC$_2$D).
	Much of it is copy-pasted from the thesis so it should not be interpreted as a publishable work.

	As the suggestive name suggestively suggests, QC$_2$D is a model where the gauge fields are instead
	elements of $SU(2)$. Thus the generators for the gauge fields are the Pauli matrices rather than the Gell-Mann
	matrices. Being a simpler theory gives it a number of computation advantages. But the main motivator for us is a
	little party trick known as the \emph{Pauli-G\"ursey} symmetry
	\begin{equation}\label{Eq:Pauli-Gursey}
		U^*=\sigma_2U\sigma_2
	\end{equation}
	which ensures that
	\begin{equation}\label{Eq:QC2D real}
		\det D=\det\left[\sigma_2 D \sigma_2\right]=\det\left[D^*\right]\in\mathbb{R}
	\end{equation}
	allowing simulations at $\mu_q\ne0$.
	
\section{Background}\label{Sec:Background}
	We are starting with an $SU(2)$ gauge theory where the gauge fields form a Lie algebra belonging to the adjoint
	representation of $SU(N)$.
	\begin{equation}\label{Eq:Lie Gauge}
		A_\mu(x)=\sum\limits_{i=1}^{N^2-1} T_iA^{(i)}_\mu(x)
	\end{equation}
	where $g$ is the coupling constant, $T_i$ form a basis for $\mathfrak{su}(N)$ and $A^{(i)}_\mu(x)$ are real-valued
	fields.

	After discretisation we can no longer have infinitesimal gauge transformations. The solution is to instead express
	the gauge fields as elements of the Lie \emph{group}
	\begin{equation}\label{Eq:Gauge field}
		U_\mu(x)=e^{ia A_\mu(x)}
	\end{equation}
	where $x$ is the lattice site and $\mu$ is the direction of the link containing the gauge field.  These gauge fields
	are written down as
	\begin{equation}\label{Eq:SU(2) Def}
		SU(2)=\left\{\begin{pmatrix} \alpha & \beta \\ -\bar{\beta} & \bar{\alpha}\end{pmatrix}; \alpha,\beta\in\mathbb{C};
		|\alpha|^2+|\beta|^2=1\right\}\text{.} 
	\end{equation}
	Using symmetries, we only need to store the top line of the matrix to carry out the simulation.  This is by no means
	a unique representation of $SU(2)$ matrices. If we consider $\alpha=a+di$ and $\beta=c+bi$ then it is not difficult
	to see that
	\begin{equation}\label{Eq:Pauli SU(2)}
		U_\mu(x)=a\mathbbm{1}+ib\sigma_x+ic\sigma_y+id\sigma_z\text{.} 
	\end{equation}
	where $\sigma_x,\sigma_y,\sigma_z$ are the Pauli matrices. This is equivalent to expressing the gauge field as a
	quaternion. Written in this form it is clear that the real component of $\alpha$ corresponds to the trace of the
	gauge link and the remaining components to the traceless components of the gauge links. The analysis code uses this
	latter form to represent the gauge fields.
	
	
	By convention, we assume that $\mu>=0$. But there is no reason that one cannot have a gauge link pointing in the
	negative direction. In fact, we can even express them in terms of a gauge link in the forwards direction
	\begin{equation}\label{Eq:Backwards Link}
	U_{-\mu}(x)=U_\mu^\dagger\left(x-\hat{\mu}\right)\text{.} 
	\end{equation}
	
	Closed paths on the lattice correspond to gauge invariant quantities, and thus observable quantities. The
	plaquette is the simplest closed loop on the lattice and is formed by a square of gauge links
	\begin{equation}\label{Eq:Plaquette}
		U_{\mu,\nu}(x)=U_\mu(x)U_\nu\left(x+\hat{\mu}\right)U_{-\mu}\left(x+\hat{\mu}+\hat{\nu}\right)U_{-\nu}\left(x+\hat{\nu}\right)
	\end{equation}
	or using \cref{Eq:Backwards Link}
	\begin{equation}\label{Eq:Plaquette Short}
		U_{\mu,\nu}(x)=U_\mu(x)U_\nu\left(x+\hat{\mu}\right)U^\dagger_{\mu}\left(x+\hat{\nu}\right)U^\dagger_{\nu}(x)\text{.} 
	\end{equation}
	Usually, we assume that $\nu>\mu$ when calculating the plaquette. This is purely a convention though.

\section{Wilson Action}
	Now that we have plaquettes, they can be used to construct the Wilson gauge action\index{Action!Wilson gauge}.
	\begin{equation}\label{Eq:Wilson Gauge Action}
		S_G[U]=\frac{\beta}{N_c}\sum\limits_{x\in\Lambda}\sum\limits_{\nu>\mu}\re{\Tr{\left[\ID-U_{\mu\nu}(x)\right]}}
	\end{equation}
	where $\Lambda$ means all sites on the lattice and
	\begin{equation}\label{Eq:Inverse Gauge Coupling}
		\beta=\frac{2N_c}{g^2}
	\end{equation}
	is the inverse gauge coupling. The $\ID$ term ensures that the action is zero when the gauge fields are set to the
	identity. In this case $A_\mu(x)=0$ so corresponds to free fields.

	If we na\"ively try to discretise the fermion component of the action, we run into the doubling problem. We get
	around this by using the Wilson Fermion action. The aim is to add a mass term to the propagator that cancels with the
	doublers yet also disappears in the continuum.
	\begin{equation}\label{Eq:Wilson term}
		\tilde{D}(p)=m+\frac{i}{a}\sum\limits_{\mu=1}^{4}\gamma_\mu\sin(p_\mu a)
		+\frac{1}{a}\sum\limits_{\mu=1}^4 \left(1-\cos\left(p_\mu a\right)\right)\text{.} 
	\end{equation}
	In the continuum limit, the final term becomes infinitely heavy so decouples from the theory. Additionally, it is $0$
	for $p=0$ but non-zero otherwise, and therefore eliminates all the double poles. Inverting the Fourier transform
	then leaves us with the Wilson-Dirac operator
	\begin{multline}\label{Eq:Wilson Dirac Op}
		D(x|y)=\\\sum\limits_{\mu=1}^4\frac{\left(1-\gamma_\mu\right)U_\mu(x)\delta_{x+\hat{\mu},y}
		-\left(1+\gamma_\mu\right)U_{-\mu}(x)\delta_{x-\hat{\mu},y}}{2a}+\left(m+\frac{4}{a}\right)\ID\delta_{x,y}\text{.} 
	\end{multline}
	If we try to take the massless limit again, we are left with a factor of $\frac{4}{a}$ in the Wilson-Dirac operator,
	and thus the discretisation errors are now $\mathscr{O}\left(a\right)$.
\section{Wilson HMC Update}
\subsection{Gauge Update}
	The gauge field updates require evaluating
	\begin{equation}\label{Eq:Gauge_Upd}
		U_\mu(x,\tau+\delta\tau)=e^{i \delta\tau \vec{\pi}_\mu(x,\tau)}U_\mu(x) 
	\end{equation}
	and since $SU(2)\cong SO(3)$ we can evaluate the exponential analytically in a similar manner as we do with complex
	numbers. Given an $SU(2)$ element written in quaternion form $A=a+\vec{b}$.
	\begin{equation}\label{Eq:SU2_Exp}
		e^{iA}=e^{ia}\left(\cos\left(|\vec{b}|\right)+i\hat{b}\sin\left(|\vec{b}|\right)\right)\text{.} 
	\end{equation}
	This follows from the Taylor series of $e^{\vec{b}}$. Even powers of $\vec{b}$ are just products of dot products so
	scalar, whereas odd powers are given by $\left|\vec{b}\right|^{n-1}\vec{b}=\left|\vec{b}\right|^n \hat{b}$. Since
	$\vec{\pi}$ doesn't have a scalar component we can drop the $e^a$ term out front. Letting
	$\pi=\pi_\mu\left(x+\frac{\delta\tau}{2}\right)$, the new gauge field is given as
	\begin{multline}\label{Eq:SU2 Gauge Update}
		U_\mu(x,\tau+\delta\tau)=\\ \delta\tau \begin{pmatrix}
		\cos\left|\vec{\pi}\right|+i \pi_z \sin\left|\vec{\pi}\right|
		& \pi_y\sin\left|\vec{\pi}\right|+i\pi_x\sin\left|\vec{\pi}\right|\\
		-\pi_y\sin\left|\vec{\pi}\right|+i\pi_x\sin\left|\vec{\pi}\right|
		& \cos\left|\vec{\pi}\right|-i\pi_z\sin\left|\vec{\pi}\right|
		\end{pmatrix} U_\mu(x,\tau)\text{.} 
	\end{multline}

	All that's left to do is an $SU(2)$ multiplication. The implementation of this is presented \verb|integrate.c|.
\subsection{Momentum updates}\label{Sub:Momentum Update}
	In order to update the conjugate momenta, we need to evaluate the force terms acting on them.
	The force is given by
	\begin{equation}\label{Eq:HMC Force}
		F\left[U,\phi\right]=\sum\limits_{i=1}^{N_c^2-1} T_i\nabla^{(i)}S\left[U\,\phi,\bar{\phi}\right]=
		\sum\limits_{i=1}^{N_c^2-1}T_i\nabla^{(i)}\left(S_G\left[U\right]+\phi^\dagger\left(M^\dagger M\right)^{-1}\phi\right)
	\end{equation}
	where $\nabla^{(i)}$ is the derivative with respect to $A$.

\subsubsection{Gauge Force}\label{Subsub:Gauge force}
	Using \cref{Eq:Gauge field} we can express the derivative of a function $f\left(U_\mu(x)\right)$ with
	respect to the gauge field $U_\mu(x)$ and generator $j$ as
	\begin{equation}\label{Eq:Gauge deriv}
			\nabla^{j}f(U_\mu(x))=\frac{\partial}{\partial\alpha}f \left[e^{i\alpha T_j}U_\mu(x)\right]_{\alpha=0}\text{.} 
	\end{equation}
	Then the gauge force is just the derivative of the plaquettes
	\begin{multline}\label{Eq:Gauge action deriv}
		T_j \nabla^{(j)}S_G\left[U\right] \\ =\frac{\beta}{N_c}
		T_j \sum\limits_{x\in\Lambda}\sum\limits_{\mu=1}^4 \Tr\left[-\frac{\partial}{\partial\alpha}
		e^{i\alpha T_j} U_\mu(x)\sum\limits_{\nu\ne\mu} \left(A_{\mu,\nu}(x)+A_{\mu,-\nu}(x)\right)\right]
		\\=\frac{\beta}{N_c} T_j \sum\limits_{x\in\Lambda}\sum\limits_{\mu=1}^4 
		\Tr\left[-iT_j U_\mu(x)\sum\limits_{\nu\ne\mu}\left( A_{\mu,\nu}(x)+A_{\mu,-\nu}(x)\right)\right]
	\end{multline}
	where $A_{\mu\nu}(x)$ is the staple in the $\mu-\nu$ direction. The plaquettes are elements of $SU\left(N_c\right)$,
	so we can use \cref{Eq:Gauge field} to express them in the form of \cref{Eq:Pauli SU(2)}, or the
	equivalent for $SU(N_c)$ if $N_c\ne2$. Since the trace of a sum is the sum of the traces, and $T_j\ID=T_j$ is
	traceless, we are left with
	\begin{equation}\label{Eq:Gauge action deriv expanded}
		T_j \nabla^{(j)}S_G\left[U\right]
		=\frac{2\beta}{N_c} \sum\limits_{x\in\Lambda}\sum\limits_{\mu=1}^4 
		T_j \Tr\left[T_j\sum\limits_{i=1}^{N_c^2-1}\left(c_i(x,\mu)T_i\right)\right]
	\end{equation}
	where $c_i(x,\mu)$ is the coefficient of generator $i$ at site $x$ in direction $\mu$, and the factor of $2$ from the
	definition of the generators. Expanding the term inside the trace, all of these products are traceless except for
	$i=j$. Taking the trace of $T_iT_j=\frac{1}{2}\delta_{ij}$ leaves us with
	\begin{equation}\label{Eq:Gauge action deriv final}
		T_j \nabla^{(j)}S_G\left[U\right]
		=\frac{\beta}{N_c} \sum\limits_{x\in\Lambda}\sum\limits_{\mu=1}^4 c_j(x,\mu)T_j\text{.} 
	\end{equation}
	Summing over all $j$, we see that this is merely the traceless component of the gauge action (up to a factor of $i$)
	and can be calculated using the same approach as pure gauge theory.

\subsubsection{Fermion Force}\label{Subsub:Fermion Force}
	The entry of the fermion matrix $M$ at index $yx$ is
	\begin{equation}\label{Eq:Fermi Matrix index}
		M_{yx}=\delta_{yx}-\kappa\sum_{\mu=\pm1}^{\pm4}\delta_{y,x+\hat{\mu}}\left(1+\gamma_\mu\right)U_{x,\mu}
	\end{equation}
	and thus
	\begin{equation}\label{Eq:Fermi Matrix dagger_index}
		(M^\dagger)_{yx}=\left(M_{yx}\right)^\dagger=
		\delta_{yx}-\kappa\sum_{\mu=\pm1}^{\pm4}U^\dagger_{y,\mu}\left(1+\gamma_\mu\right)\delta_{y,x+\hat{\mu}}\text{.} 
	\end{equation}

	If we now look at the fermionic contribution to the force we can apply the product rule and obtain
	\begin{equation}\label{Eq:M_dag M deriv def}
		\nabla^{(j)}\left(M^\dagger M\right)^{-1}=
		\left(M^\dagger M\right)^{-1}\left[\left(\nabla^{(j)}
		M^\dagger\right)M+M^\dagger\left(\nabla^{(j)}M\right)\right]\left(M^\dagger M\right)^{-1}
	\end{equation}
	Using \cref{Eq:Gauge deriv}, the derivative of the fermion matrix at site $x$ in direction $\mu$ for colour
	$j$ is given by
	\begin{align}\label{Eq:Fermi Matrix Deriv}
		D_{x\mu j}M_{yx}&= -i\kappa\delta_{y,x+\mu}\left(1+\gamma_\mu\right)T_j U_{x,\mu}\nonumber\\
		D_{x\mu j}M^\dagger_{yx}&= i\kappa U^\dagger_{y,\mu}T_j\left(1+\gamma_\mu\right)\delta_{y,x+\hat{\mu}}\text{.} 
	\end{align}
	Applying these derivative terms to \cref{Eq:M_dag M deriv def} gives
	\begin{multline}\label{Eq:Fermi Prod Rule 1}
		\left(M_{zy}\right)^\dagger\left(D_{x\mu j}M_{zx}\right)=\\
		\left[\delta_{yz}-\kappa\sum\limits_{\nu=\pm1}^{\pm4}
		U^\dagger_{y,\nu}\left(1+\gamma_\nu\right)\delta_{y,x+\hat{\nu}}\right]
		\left[-i\kappa\delta_{z,x+\hat{\mu}}\left(1+\gamma_\mu\right)T_j U_{x,\mu}\right]=\\
		-i\kappa\delta_{yz}\delta_{z,x+\hat{\mu}}\left(1+\gamma_\mu\right)T_j U_{x,\mu}+
		i\kappa^2\sum\limits_{\nu=\pm1}^{\pm4}U^\dagger_{y\nu}\left(1+\gamma_\nu+\gamma_\mu+\gamma_\nu\gamma_\mu\right)
		\delta_{y,z+\hat{\nu}}\delta_{z,x+\hat{\mu}}T_j U_{x\mu}\text{.} 
	\end{multline}
	and
	\begin{multline}\label{Eq:Fermi Prod Rule 2}
		\left(D_{x\mu j}\left(M_{zy}\right)^\dagger\right)M_{zx}=\\
		\left[i\kappa U^\dagger_{y,\mu}T_j\left(1+\gamma_\mu\right)\delta_{y,z+\hat{\mu}}\right]
		\left[\delta_{zx}-\kappa\sum\limits_{\nu=\pm1}^{\pm4}
		\delta_{z,x+\hat{\nu}}\left(1+\gamma_\nu\right)U^\dagger_{x,\nu}\right]=\\
		i\kappa U^\dagger_{y,\mu}T_j\left(1+\gamma_\mu\right)\delta_{y,z+\hat{\mu}}\delta_{zx} 
		-i\kappa^2\sum\limits_{\nu=\pm1}^{\pm4}U^\dagger_{y\mu}T_j\left(1+\gamma_\mu+\gamma_\nu+\gamma_\mu\gamma_\nu\right)
		\delta_{y,z+\hat{\mu}}\delta_{z,x+\hat{\nu}}U_{x\nu}\text{.} 
	\end{multline}

	The $\kappa^2$ terms in both \cref{Eq:Fermi Prod Rule 1} and \cref{Eq:Fermi Prod Rule 2} are non-zero for
	$y=z+\hat{\mu}=x+\hat{\mu}+\hat{\nu}$. As we sum over $\nu$ and $\mu$ we can swap the $\mu$ and $\nu$ indices and
	these terms cancel each other out. The $\kappa$ terms are non-zero for $y=z=x+\hat{\mu}$ in \cref{Eq:Fermi Prod
	Rule 1} and $y=z=x-\hat{\mu}$ in \cref{Eq:Fermi Prod Rule 2} respectively. This finally leaves us with an expression
	for the derivative at site $x$ in direction $\mu$ for colour index $j$
	\begin{multline}\label{Eq:M_dag M deriv}
		D_{x\mu j}\left(M^\dagger M\right)^{-1}=\\\left(M^\dagger
		M\right)^{-1}_{xz}\left(
		-i\kappa\left(1+\gamma_\mu\right)T_j U_{x,\mu}+
		i\kappa U^\dagger_{x-\hat{\mu},\mu}T_j\left(1+\gamma_\mu\right)
		\right)\left(M^\dagger M\right)^{-1}_{zx}\text{.} 
	\end{multline}
\section{Clover Fermions}
	The Clover improved Wilson action can be given as
	\begin{equation}
		S_\text{I}=S_\text{Wilson}+S_\text{Clov}=S_\text{Wilson}
		-\frac{c_\text{SW}}{2} \kappa a^5 \sum\limits_{n\in\Lambda}\sum\limits_{\mu<\nu}
		\bar{\psi}(n)\frac{1}{2}\sigma_{\mu\nu}\hat{F}_{\mu\nu}(n)\psi(n)
	\end{equation}
	where $c_\text{SW}$ is the clover or Sheiholeslami-Wohlert coefficient, $\sigma_{\mu\nu}$ is the commutator of the
	Euclidean Dirac matrices in the chiral representation and the clover term at site $x$ in the $\mu-\nu$ direction is
	given by
	\begin{equation}\label{Eq:Clover}
		F_{\mu\nu}(x)=\frac{-i}{8a^2}\left(Q_{\mu\nu}(x)-Q_{\nu\mu}(x)\right)
	\end{equation}
	where (at least in our implementation, it is not unique)
	\begin{figure}
		\centring
		\begin{adjustbox}{width=\linewidth}
			\tikzfig{Tikz/Clover}
		\end{adjustbox}
		\captionsetup{width=\linewidth}
		\caption{A half-clover $Q_{\mu\nu}$. The full clover term consists of this minus its hermitian conjugate}
		\label{Tikz:Clover}
	\end{figure}
	\begin{equation}\label{Eq:Half Clover}
		Q_{\mu\nu}=U_{\mu,\nu}(x)+U_{\nu,-\mu}(x)+U_{-\nu,\mu}(x)+U_{-\mu,-\nu}(x)
	\end{equation}
	and $U_{\mu,\nu}$ is the plaquette given in \cref{Eq:Plaquette} (see \cref{Tikz:Clover}). Since the
	$Q_{\mu\nu}$ terms are the sum of four plaquettes they should have the same form as $U_\mu$ given in \cref{Eq:SU(2)
	Def} (but are \emph{not} an element of $SU(2)$), from here it follows that
	\begin{equation}\label{Eq:Diff_form}
		Q_{\mu\nu}-Q_{\nu\mu}=\begin{pmatrix}
		\alpha-\bar{\alpha} & \beta+\beta \\
		-\bar{\beta}-\bar{\beta} & \bar{\alpha}-\alpha \end{pmatrix}=
		\begin{pmatrix} 2di & 2c+2bi \\ -2c+2bi & -2di \end{pmatrix}
	\end{equation}
	which is anti-Hermitian. Multiplying out front by $-i$ gives us 
	\begin{equation}\label{Eq:Clov_Form}
		F_{\mu\nu}(x)=\frac{1}{8a^2}
		\begin{pmatrix} 2d & 2b-2ci \\ 2b+2ci & -2d \end{pmatrix}\text{.}
	\end{equation}
	which is hermitian and traceless (i.e. of the same form as an element of the Lie Algebra).
\section{Clover Force Update}
	For the clover force update with respect to a link $U_\mu$ is similar to the procedure described in 
	\cref{Subsub:Fermion Force}. Only the leaves containing $U_\mu(x)$ or $U^\dagger_{\mu}\left(x+\hat{\mu}\right)$
	contribute. It does have the additional complication that there are six clovers containing the gauge link between $x$
	and $x+\hat{\mu}$ or the conjugate link between $x+\hat{\mu}$ and $x$.
	\begin{enumerate}
		\item The clover at the site $x$ itself contains the link in the $+\mu,+\nu$ leaf and the $+\mu,-\nu$ leaf.
		\item The clover at the site $x+\hat{\mu}$ contains the link in the $-\mu,+\nu$ leaf and the $-\mu,-\nu$ leaf.
		\item The clover at the site $x+\hat{\nu}$ contains the link in the $+\mu,-\nu$ leaf.
		\item The clover at the site $x-\hat{\nu}$ contains the link in the $+\mu,+\nu$ leaf.
		\item The clover at the site $x+\hat{\mu}+\hat{\nu}$ contains the link in the $-\mu,-\nu$ leaf.
		\item The clover at the site $x+\hat{\mu}-\hat{\nu}$ contains the link in the $-\mu,+\nu$ leaf.
	\end{enumerate}
	\begin{figure}
		\centring
		\begin{adjustbox}{width=\linewidth}
			\tikzfig{Tikz/Clover_Force_All}
		\end{adjustbox}
		\captionsetup{width=\linewidth}
		\caption{Leaves containing the link $U_\mu(x)$ or $U^\dagger_{\mu}\left(x+\hat{\mu}\right)$ centred on the
		rectangles Note that all the clovers drawn in the diagram are $Q_{\mu\nu}$. The contribution from $Q_{\nu\mu}$ is
		not shown.}
		\label{Tikz:Clover_Force}
	\end{figure}

	Unlike the Wilson case where the generator multiplies the plaquette from the left, the location of the generator in
	the clover force depends on the location of the link in each clover. For the $Q_{\mu\nu}$ leaves seen in 
	\cref{Tikz:Clover_Force,Tikz:Clover_Force_Comp}.
	\begin{enumerate}
		\item The clover at the site $x$ itself has the generator multiplying its top-right leaf from the left (w.r.t. $U_\mu(x)$)
				and bottom right leaf from the right (w.r.t. $U^\dagger_\mu\left(x+\hat{\mu}\right)$).
		\item The clover at the site $x+\hat{\mu}$ has the generator appearing between the third and fourth link of its
				top left leaf ($\nu,-\mu$) (w.r.t. $U_\mu(x)$) and between the first and second link of its bottom left leaf
				($-\mu,-\nu$) (w.r.t. $U^\dagger_\mu\left(x+\hat{\mu}\right)$).
		\item The clover at the site $x+\hat{\nu}$ has the generator appearing between the first and second link (w.r.t. $U_\mu(x)$)
		\item The clover at the site $x-\hat{\nu}$ has the generator appearing between the third and fourth link
				(w.r.t. $U^\dagger_\mu\left(x+\hat{\mu}\right)$).
		\item The clover at the site $x+\hat{\mu}+\hat{\nu}$ has the generator appearing between second and third link (w.r.t. $U_\mu(x)$).
		\item The clover at the site $x+\hat{\mu}-\hat{\nu}$ has the generator appearing between second and third link
				(w.r.t. $U^\dagger_\mu\left(x+\hat{\mu}\right)$).
	\end{enumerate}
	\begin{figure}
		\centring
		\begin{adjustbox}{width=\linewidth}
			\tikzfig{Tikz/Clover_Force_Comp}
		\end{adjustbox}
		\captionsetup{width=\linewidth}
		\caption{The components of the clover force with respect to $U_\mu(x)$ (the link with arrow pointing \emph{away}
		from the centre of the respective clover.) Negative terms are from $Q_{\nu\mu}(x)$ terms. Diagonal lines indicate
		where the generator appears in the derivative.}
		\label{Tikz:Clover_Force_Comp}
	\end{figure}

	Looking at the $F_{\mu\nu}(x)$ and the derivatives with respect to $U_\mu(x)$, we see that the derivative from each
	generator takes the form of an $SU(2)$ group element
	\begin{align}
		\frac{\partial U_{\mu\nu}(x)}{\partial U_\mu(x)}&=
		i\sigma_x U_\mu(x)U_\nu\left(x+\hat{\mu}\right)U^\dagger_\mu\left(x+\hat{\nu}\right)U^\dagger_\nu(x)\nonumber\\
		&=\begin{pmatrix} -i\bar{\beta} & i\bar{\alpha} \\i\alpha & i\beta\end{pmatrix}=
		\begin{pmatrix} -b-ci & d+ai \\ -d+ai & -b+ci\end{pmatrix}\label{Eq:Sig_x_deriv}\\
		\frac{\partial U_{\mu\nu}(x)}{\partial U_\mu(x)}&=
		i\sigma_y U_\mu(x)U_\nu\left(x+\hat{\mu}\right)U^\dagger_\mu\left(x+\hat{\nu}\right)U^\dagger_\nu(x)\nonumber\\
		&=\begin{pmatrix} \bar{\beta} & -\bar{\alpha} \\\alpha & \beta\end{pmatrix}=
		\begin{pmatrix} c-bi & -a+di \\ a+di & c+bi\end{pmatrix}\label{Eq:Sig_y_deriv}\\
		\frac{\partial U_{\mu\nu}(x)}{\partial U_\mu(x)}&=
		i\sigma_x U_\mu(x)U_\nu\left(x+\hat{\mu}\right)U^\dagger_\mu\left(x+\hat{\nu}\right)U^\dagger_\nu(x)\nonumber\\
		&=\begin{pmatrix} i\alpha & i\beta \\i\bar{\beta} & -i\bar{\alpha}\end{pmatrix}=
		\begin{pmatrix} -d+ai & -b+ci \\ b+ci & -d-ai\end{pmatrix}\label{Eq:Sig_z_deriv}\text{.}
	\end{align}
	Let's look at a specific example. Taking $\lambda_j=\sigma_x$ as the generator and using the clover at the site, we
	see that the contribution from the $Q_{\nu\mu}=Q^\dagger_{\mu\nu}$ leaf
	\begin{equation}\label{Eq:Site_Leaf_Conj_Deriv}
		\frac{\partial U^\dagger_{\mu\nu}(x)}{\partial U^\dagger_\mu\left(x+\hat{\mu}\right)}=
		-i U_\nu(x)U_\mu\left(x+\hat{\nu}\right)U^\dagger_\nu\left(x+\hat{\mu}\right)U^\dagger_\mu(x)\lambda_j=
		\left(\frac{\partial U_{\mu\nu}(x)}{\partial U_\mu(x)}\right)^\dagger\text{.}
	\end{equation}
	\Cref{Eq:Site_Leaf_Conj_Deriv} is useful as it means we do not have to worry about calculating and keeping track of
	$U^\dagger_{\mu\nu}(x)$ at the same time as $U_{\mu\nu}(x)$. Just work with the undaggered term instead. We can apply
	this symmetry to the $U_{-\nu,\mu}(x)$ leaf to obtain $\frac{\partial U_{\mu,-\nu}(x)}{\partial U_\mu(x)}$
	\begin{equation}\label{Eq:Deriv_Symm}
		\frac{\partial U_{-\nu,\mu}(x)}{\partial U^\dagger_\mu(x)}=
		\left(\frac{\partial U^\dagger_{-\nu,\mu}(x)}{\partial U_\mu(x)}\right)^\dagger\Leftrightarrow
		\frac{\partial U_{\mu,-\nu}(x)}{\partial U_\mu(x)}=
		\frac{\partial U^\dagger_{-\nu,\mu}(x)}{\partial U_\mu(x)}=
		\left(\frac{\partial U_{-\nu,\mu}(x)}{\partial U^\dagger_\mu(x)}\right)^\dagger\text{.}
	\end{equation}

	Using this symmetry and getting the difference yields
	\begin{multline}\label{Eq:Clov_Force_1}
		\frac{\partial F_{\mu\nu}(x)}{\partial U_\mu(x)}=\frac{-i}{8a}\left(\frac{\partial Q_{\mu\nu}(x)}
		{\partial U_\mu(x)}-\frac{\partial Q_{\nu\mu}(x)}{\partial U_\mu(x)}\right)\\
		=\frac{-i}{8a}\left(\frac{\partial U_{\mu\nu}(x)}{\partial U_\mu(x)}
		-\frac{\partial U_{\mu,-\nu}(x)}{\partial U_\mu(x)}\right)
		=\frac{-i}{8a}\left(\frac{\partial U_{\mu\nu}(x)}{\partial U_\mu(x)}
		-\left(\frac{\partial U_{-\nu,\mu}(x)}{\partial U^\dagger_\mu\left(x+\hat{\mu}\right)}\right)^\dagger\right)
	\end{multline}
	as the total contribution to the derivative of the clover centred at $x$ with respect to $U_\mu(x)$. Repeating for
	the other contributions
	\begin{align}
		%x+mu
		\frac{\partial F_{\mu\nu}\left(x+\hat{\mu}\right)}{\partial U_\mu(x)}&=
		\frac{-i}{8a}\left(\frac{\partial Q_{\mu\nu}\left(x+\hat{\mu}\right)}{\partial U_\mu(x)}-\frac{\partial
		Q_{\nu\mu}\left(x+\hat{\mu}\right)}{\partial U_\mu(x)}\right)\nonumber\\
		&=\frac{-i}{8a}\left(\frac{\partial U_{\nu,-\mu}\left(x+\hat{\mu}\right)}{\partial U_\mu(x)}
		-\frac{\partial U_{-\nu,-\mu}\left(x+\hat{\mu}\right)}{\partial U_\mu(x)}\right)\nonumber\\
		&=\frac{-i}{8a}\left(\frac{\partial U_{\nu,-\mu}\left(x+\hat{\mu}\right)}{\partial U_\mu(x)}
		-\left(\frac{\partial U_{-\mu,-\nu}\left(x+\hat{\mu}\right)}{\partial U^\dagger_\mu(x)}\right)^\dagger\right)\label{Eq:Clov_Force_2}\\
		% x+nu
		\frac{\partial F_{\mu\nu}\left(x+\hat{\nu}\right)}{\partial U_\mu(x)}&=
			\frac{-i}{8a}\frac{\partial Q_{\mu\nu}\left(x+\hat{\nu}\right)}{\partial U_\mu(x)}=
			\frac{-i}{8a}\frac{\partial U_{-\nu,\mu}\left(x+\hat{\nu}\right)}{\partial U_\mu(x)}\label{Eq:Clov_Force_3}\\
		% x-nu
		\frac{\partial F_{\mu\nu}\left(x-\hat{\nu}\right)}{\partial U_\mu(x)}&=
			\frac{i}{8a}\frac{\partial Q_{\nu\mu}\left(x-\hat{\nu}\right)}{\partial U_\mu(x)}\nonumber\\
			&=\frac{i}{8a}\frac{\partial U_{\nu,\mu}\left(x-\hat{\nu}\right)}{\partial U_\mu(x)}=
			\frac{i}{8a}\left(\frac{\partial U_{\mu,\nu}\left(x-\hat{\nu}\right)}
			{\partial U^\dagger_\mu(x)}\right)^\dagger\label{Eq:Clov_Force_4}\\
		% x+mu+nu
		\frac{\partial F_{\mu\nu}\left(x+\hat{\mu}+\hat{\nu}\right)}{\partial U_\mu(x)}&=
			\frac{-i}{8a}\frac{\partial Q_{\mu\nu}\left(x+\hat{\mu}+\hat{\nu}\right)}{\partial U_\mu(x)}=
			\frac{-i}{8a}\frac{\partial U_{-\mu,-\nu}\left(x+\hat{\mu}+\hat{\nu}\right)}{\partial U_\mu(x)}\label{Eq:Clov_Force_5}\\
		% x+mu-nu
		\frac{\partial F_{\mu\nu}\left(x+\hat{\mu}-\hat{\nu}\right)}{\partial U_\mu(x)}&=
			\frac{i}{8a}\frac{\partial Q_{\nu\mu}\left(x+\hat{\mu}-\hat{\nu}\right)}{\partial U_\mu(x)}\nonumber\\
			&=\frac{i}{8a}\frac{\partial U_{-\mu,\nu}\left(x+\hat{\mu}-\hat{\nu}\right)}{\partial U_\mu(x)}
			=\frac{i}{8a}\left(\frac{\partial U_{\nu,-\nu}\left(x+\hat{\mu}-\hat{\nu}\right)}
			{\partial U^\dagger_\mu(x)}\right)^\dagger\label{Eq:Clov_Force_6}\text{.}
	\end{align}

	Next, we need to consider the contribution from the derivative with respect to
	$U^\dagger_\mu\left(x+\hat{\mu}\right)$. The procedure for this is similar to that for 
	\cref{Eq:Clov_Force_1,Eq:Clov_Force_2,Eq:Clov_Force_3,Eq:Clov_Force_4,Eq:Clov_Force_5,Eq:Clov_Force_6}.
	Focusing only on $F_{\mu\nu}(x)$ and using \cref{Eq:Deriv_Symm}
	\begin{multline}\label{Eq:Clover_deriv_dag}
		\frac{\partial F_{\mu\nu}(x)}{\partial U^\dagger_\mu\left(x+\hat{\mu}\right)}=
		\frac{-i}{8a}\left(\frac{\partial Q_{\mu\nu}(x)} {\partial U^\dagger_\mu\left(x+\hat{\mu}\right)}
		-\frac{\partial Q_{\nu\mu}(x)}{\partial U^\dagger_\mu\left(x+\hat{\mu}\right)}\right)\\
		=\frac{-i}{8a}\left(\frac{\partial U_{\nu,-\mu}(x)}{\partial U^\dagger_\mu\left(x+\hat{\mu}\right)}-
			\frac{\partial U_{\nu,\mu}(x)}{\partial U^\dagger_\mu\left(x+\hat{\mu}\right)}\right)
		=\frac{-i}{8a}\left(\frac{\partial U_{\nu,-\mu}(x)}{\partial U^\dagger_\mu\left(x+\hat{\mu}\right)}-
			\left(\frac{\partial U_{\mu,\nu}(x)}{\partial U_\mu(x)}\right)^\dagger\right)\text{.}
	\end{multline}
	Comparing to \cref{Eq:Clov_Force_1}, we can quickly see that
	\begin{equation}\label{Eq:NoDag_vs_Dag}
		\frac{\partial F_{\mu\nu}(x)}{\partial U^\dagger_\mu\left(x+\hat{\mu}\right)}=
		\left(\frac{\partial F_{\mu\nu}(x)}{\partial U_\mu(x)}\right)^\dagger
	\end{equation}
	and thus from \cref{Eq:Sig_x_deriv} the overall form of the derivatives in the $\hat{\mu}$ direction is
	\begin{equation}
		\frac{\partial F_{\mu\nu}(x)}{\partial U_\mu(x)}+
		\frac{\partial F_{\mu\nu}(x)}{\partial U^\dagger_\mu\left(x+\hat{\mu}\right)}=
		\frac{-i}{8a}\begin{pmatrix} -i\bar{\beta} +i\beta& i\bar{\alpha}-i\bar{\alpha} \\
		i\alpha-i\alpha & i\beta-i\bar{\beta}\end{pmatrix}=
		\frac{1}{8a}\begin{pmatrix} 2bi & 0 \\ 0 & 2bi\end{pmatrix}
	\end{equation}
	which is anti-hermitian. Comparing this with \cref{Eq:Pauli SU(2)}, we see that the force from the $\sigma_x$
	generator only consists of the $\sigma_x$ components of the leaves involved in the force. Using a similar approach we
	can also show that the $\sigma_y$ and $\sigma_z$ components behave in the same manner.

	Unfortunately, there is no symmetry relating $\frac{\partial F_{\mu\nu}(x)}{\partial U_\mu(x)}$ and
	$\frac{\partial F_{\nu\mu}(x)}{\partial U_\nu(x)}$ so these must be calculated explicitly. Thankfully this is just a
	matter of calling the clover force routine with $\mu$ and $\nu$ swapped. 
\end{document}
	Lastly, we multiply this by $\frac{-i}{8}$ to get
	\begin{equation}\label{Eq:Clov_Force_Deriv}
		\frac{\partial F_{\mu\nu}}{\partial	U_\mu}=\frac{1}{8}
		\begin{pmatrix} -2c & -2a-2di \\ 2a+2di & 2c\end{pmatrix}
	\end{equation}
	which if I am reading this correctly is neither hermitian nor anti-hermitian. This implies that there is a mistake
	somewhere. I am confident that all the other terms will have a similar form, but will not check here to save time.
